\documentclass[11pt,a4paper]{article}

\usepackage{amsmath,amssymb,amsthm,mathtools}
\usepackage[utf8]{inputenc}
\usepackage{mathrsfs}
\usepackage[margin=1in]{geometry}
\usepackage{hyperref}
\usepackage{enumitem}

% Theorem environments
\theoremstyle{plain}
\newtheorem{theorem}{Theorem}[section]
\newtheorem{proposition}[theorem]{Proposition}
\newtheorem{lemma}[theorem]{Lemma}
\newtheorem{corollary}[theorem]{Corollary}

\theoremstyle{definition}
\newtheorem{definition}[theorem]{Definition}
\newtheorem{example}[theorem]{Example}

\theoremstyle{remark}
\newtheorem{remark}[theorem]{Remark}

% Operators
\newcommand{\Id}{\mathrm{Id}}
\newcommand{\E}{\mathbb{E}}
\newcommand{\R}{\mathbb{R}}
\newcommand{\N}{\mathbb{N}}
\newcommand{\im}{\mathrm{im}}
\DeclareMathOperator{\PV}{P.V.}

\title{Two-Parameter Clark--Ocone Formulas in Banach Energy Spaces\\
for Stable L\'evy Processes}

\author{Ramiro Fontes\\
Quijotic Research\\
\texttt{ramirofontes@gmail.com}}

\date{February 2026}

\begin{document}

\maketitle

\begin{abstract}
We extend the operator factorization framework
of~\cite{Fontes2026a,Fontes2026d} from one-parameter to
two-parameter energy spaces for symmetric $\gamma$-stable
L\'evy processes ($\gamma \in (1,2)$).

In the one-parameter framework~\cite{Fontes2026d}, the energy
space $H_L = L^p([0,T]\times\Omega,\,dt\otimes dP)$ supports
the stochastic integral $\delta_L(u) = \int u_t\,dL_t$ and
the adjoint derivative $D_L = \delta_L^*$.  The Banach It\^o
formula~\cite{Fontes2026f} decomposes $\phi(L_T)$ into a
$\delta_L$-integral plus a Leibniz defect
$\Gamma_L^{\phi}(T)$.  However, the defect itself is not a
$\delta_L$-integral: it lives outside the range of the
one-parameter divergence.  Consequently, the one-parameter
factorization $(\Id - \E) = \delta_L \circ \iota_L^{-1} \circ D_L$
cannot yield a Clark--Ocone representation for arbitrary
stable functionals.

We resolve this by introducing the \emph{two-parameter energy
space} $\widetilde{H}_L = L^p([0,T]\times\R\times\Omega,\,
dt\otimes\nu\otimes dP)$, where $\nu$ is the L\'evy measure,
with divergence
$\widetilde{\delta}_L(h) = \int_0^T\!\int_{\R}h(s,z)\,
\widetilde{N}(ds,dz)$ (integration against the compensated
Poisson random measure) and derivative
$\widetilde{D}_L = \widetilde{\delta}_L^*$.  The
two-parameter BDG equivalence
$\|\widetilde{\delta}_L(h)\|_{L^p} \asymp
\|h\|_{\widetilde{H}_L}$
gives the recovery map $\widetilde{\iota}_L^{-1}$, and the
factorization
$(\Id - \E) = \widetilde{\delta}_L \circ
\widetilde{\iota}_L^{-1} \circ \widetilde{D}_L$
yields an $L^p$ Clark--Ocone formula:
every centered $F \in L^p(\Omega)$ has a representation
$F - \E[F] = \widetilde{\delta}_L(h)$ for a unique
$h \in \widetilde{H}_L$.

The one-parameter and two-parameter frameworks occupy
complementary regimes.  The $\delta_L$-integrals have
integrands $u(t)\cdot z$ that are \emph{$\nu$-singular}
($\int|z|^p\nu(dz) = \infty$): they converge by cancellation
in the compensation, not by integrability.
The Leibniz defect $\Gamma_L^{\phi}$ has integrand
$\phi(L_{s-}+z)-\phi(L_{s-})-\phi'(L_{s-})z = O(z^2)$
near $z = 0$, which is \emph{$\nu$-regular}
($\int_{|z|\leq 1}|z|^{2p}\nu(dz) < \infty$ for
$p > \gamma/2$) and lies naturally in
$\widetilde{H}_L$.  The defect is thus the explicitly
representable part of $\phi(L_T)$ in the two-parameter space;
the stochastic integral part has only an abstract
$\widetilde{H}_L$-representative.

As applications, we derive regularity conditions on stable
functionals in terms of the $z$-decay of the defect
integrand, connecting the operator framework to Besov-type
regularity scales for the fractional Laplacian.
\end{abstract}

\medskip
\noindent\textbf{MSC 2020:} 60H07, 60H05, 60G52, 60G51, 46B28.\\
\textbf{Keywords:} Clark--Ocone formula; Banach energy spaces;
operator factorization; stable L\'evy processes;
Leibniz defect; compensated Poisson measure;
two-parameter representation.


%======================================================================
\section{Introduction}\label{sec:intro}
%======================================================================

\subsection{Background and motivation}

The operator factorization framework developed
in~\cite{Fontes2026a,Fontes2026b,Fontes2026c} expresses
stochastic calculus through the identity
\begin{equation}\label{eq:factorization}
  (\Id - \E) = \delta \circ \iota^{-1} \circ r \circ D,
\end{equation}
where $D$ is the operator--covariant derivative, $\delta = D^*$
is the divergence (stochastic integral), $r$ is the restriction
to predictable processes, and $\iota^{-1}$ is the recovery map.
In the Hilbert framework ($L^2$-energy spaces, Brownian filtrations),
this factorization yields the Clark--Ocone representation: every
centered $F \in L^2(\Omega)$ has a unique stochastic integral
representation $F - \E[F] = \int_0^T u_t\,dB_t$.

The extension to Banach energy spaces~\cite{Fontes2026d}
replaced $L^2$ by $L^p$ with $p < \gamma < 2$ for symmetric
$\gamma$-stable L\'evy processes.  The Banach
It\^o formula~\cite{Fontes2026f} was derived by time
discretization, yielding for $\phi \in C_b^3(\R)$:
\begin{equation}\label{eq:banach-ito}
  \phi(L_T) - \phi(0)
  = \underbrace{\int_0^T \phi'(L_{t-})\,dL_t}_{%
  \text{$\delta_L$-integral}}
  + \underbrace{\Gamma_L^{\phi}(T)}_{%
  \text{Leibniz defect}},
\end{equation}
where $\Gamma_L^{\phi}(T) = \sum_{s \leq T}
[\phi(L_s) - \phi(L_{s-}) - \phi'(L_{s-})\Delta L_s]$ is the
chain-rule defect.

The It\^o formula is a representation of $\phi(L_T)$ as a
$\delta_L$-integral plus a correction.  It is \emph{not} a
Clark--Ocone formula, because the correction
$\Gamma_L^{\phi}(T)$ is not itself a $\delta_L$-integral.


\subsection{The representability obstruction}\label{sec:obstruction}

The obstruction is structural.  Write
$\delta_L(u) = \int_0^T u_t\,dL_t$ for the one-parameter
stochastic integral.  Via the L\'evy--It\^o decomposition,
\[
  \delta_L(u)
  = \int_0^T\!\!\int_{\R} u_t\,z\,\widetilde{N}(dt,dz),
\]
where $\widetilde{N}$ is the compensated Poisson random measure.
The integrand $u_t \cdot z$ is \emph{linear in~$z$}.  Thus the
range of $\delta_L$ consists of $\widetilde{N}$-integrals with
integrands that are linear in the jump size.

The Leibniz defect has the Poisson integral representation
\begin{equation}\label{eq:defect-poisson}
  \Gamma_L^{\phi}(T)
  = \int_0^T\!\!\int_{\R}
  \bigl[\phi(L_{s-}+z) - \phi(L_{s-}) - \phi'(L_{s-})z\bigr]\,
  N(ds,dz).
\end{equation}
The integrand is the Taylor remainder
$\phi(x+z) - \phi(x) - \phi'(x)z$, which is
\emph{quadratic in~$z$} to leading order.  This integrand
is not of the form $u(s)\cdot z$, so $\Gamma_L^{\phi}(T)$ does
not lie in the range of $\delta_L$.

In the Hilbert case ($\gamma \to 2$, Brownian motion),
$\Gamma_L^{\phi} = 0$: the defect vanishes, the
one-parameter integral captures everything, and the
factorization~\eqref{eq:factorization} yields Clark--Ocone.
For $\gamma < 2$, the defect is nonzero, and the one-parameter
factorization cannot produce a full representation theorem.

\begin{remark}[Martingale representation for L\'evy processes]
\label{rem:levy-rep}
This is consistent with the classical theory: the filtration of
a L\'evy process does \emph{not} have the predictable
representation property with respect to $dL$-integrals alone.
The full martingale representation requires integration against
the compensated Poisson measure $\widetilde{N}(ds,dz)$ with
integrands depending on both $s$ and $z$
(see e.g.~\cite{Applebaum2009}, Theorem~5.3.5).
\end{remark}


\subsection{Main results}

We resolve the representability obstruction by enlarging the
energy space from one parameter ($t$) to two parameters ($t,z$).

\medskip\noindent
\textbf{Theorem~A} (Two-parameter factorization;
Theorem~\ref{thm:two-param-factorization}).
\emph{The identity
$(\Id - \E) = \widetilde{\delta}_L \circ
\widetilde{\iota}_L^{-1} \circ \widetilde{D}_L$
holds on $L^p(\Omega) \cap L^q(\Omega)$, where
$\widetilde{\delta}_L$, $\widetilde{D}_L$, and
$\widetilde{\iota}_L^{-1}$ are the two-parameter divergence,
derivative, and recovery map defined in
Section~\ref{sec:two-param}.}

\medskip\noindent
\textbf{Theorem~B} ($L^p$ Clark--Ocone;
Theorem~\ref{thm:clark-ocone}).
\emph{Every centered $F \in L^p(\Omega) \cap L^q(\Omega)$ has a
unique representation}
\[
  F - \E[F]
  = \int_0^T\!\!\int_{\R} h_F(s,z)\,\widetilde{N}(ds,dz)
\]
\emph{with $h_F \in \widetilde{H}_L
= L^p([0,T]\times\R\times\Omega,\,
\mathcal{P}\otimes\mathcal{B}(\R),\,dt\otimes\nu\otimes dP)$,
and the integrand satisfies
$h_F = \widetilde{\iota}_L^{-1}(\widetilde{D}_L F)$.}

\medskip\noindent
\textbf{Theorem~C} (Two-regime decomposition;
Theorem~\ref{thm:defect-decomposition}).
\emph{For $F = \phi(L_T)$ with $\phi \in C_b^3(\R)$, the
martingale part of the defect $\Gamma_L^{\phi}(T)$ is
$M_T^{(1)} = \widetilde{\delta}_L(g_\phi)$ with natural
integrand
$g_\phi(s,z) = \phi(L_{s-}+z) - \phi(L_{s-}) - \phi'(L_{s-})z
= O(z^2)$, which is $\nu$-regular
($g_\phi \in \widetilde{H}_L$).  The stochastic integral
$\delta_L(\phi'(L_{\cdot-}))$ also has a
$\widetilde{H}_L$-representative, but it is non-explicit:
the natural integrand $\phi'(L_{t-})z$ is $\nu$-singular
($\int|z|^p\nu(dz) = \infty$) and does not belong to
$\widetilde{H}_L$.}


\subsection{Significance}

These results complete the operator factorization program:

\begin{center}
\begin{tabular}{lll}
\textbf{Setting} & \textbf{Energy space} & \textbf{Representation} \\
\hline
Hilbert ($\gamma = 2$, BM)
  & $H_B = L^2(dt \otimes dP)$
  & 1-parameter Clark--Ocone \\
Banach, 1-param ($\gamma < 2$)
  & $H_L = L^p(dt \otimes dP)$
  & It\^o formula + defect \\
Banach, 2-param ($\gamma < 2$)
  & $\widetilde{H}_L = L^p(dt \otimes \nu \otimes dP)$
  & Full Clark--Ocone
\end{tabular}
\end{center}

The Leibniz defect $\Gamma_L$ acquires a dual interpretation:
\begin{enumerate}[label=(\roman*)]
\item In the It\^o formula: the \emph{correction term}
  (the chain-rule defect added to $\delta_L(\phi'(L_{\cdot-}))$).
\item In the representation theory: the \emph{$\nu$-regular
  component} of $F$ (the part with an explicit
  $\widetilde{H}_L$-integrand $g_\phi = O(z^2)$).
\item In the one-parameter framework: the \emph{obstruction}
  to $dL$-representability (the part of $F$ not in
  $\im(\delta_L)$).
\end{enumerate}
In the Brownian case ($\nu = 0$), the defect vanishes,
the obstruction disappears, and the one-parameter
representation suffices.  The transition from
$\gamma < 2$ to Brownian motion is singular:
the jump structure that necessitates the two-parameter
framework vanishes entirely.


\subsection{Organization}

Section~\ref{sec:one-param-recap} recalls the one-parameter
framework.  Section~\ref{sec:two-param} constructs the
two-parameter energy space, divergence, and derivative.
Section~\ref{sec:BDG} establishes the two-parameter BDG
equivalence and the recovery map.
Section~\ref{sec:factorization} proves the two-parameter
factorization and Clark--Ocone formula (Theorems~A and~B).
Section~\ref{sec:decomposition} proves the two-regime
decomposition (Theorem~C) and analyzes the relationship
between the one-parameter and two-parameter frameworks.
Section~\ref{sec:regularity} develops regularity implications.
Section~\ref{sec:discussion} discusses connections and open
problems.


%======================================================================
\section{The One-Parameter Framework (Recap)}
\label{sec:one-param-recap}
%======================================================================

We recall the setup from~\cite{Fontes2026d,Fontes2026f}.
Let $L = (L_t)_{t \in [0,T]}$ be a symmetric $\gamma$-stable
L\'evy process with $\gamma \in (1,2)$ on
$(\Omega, \mathcal{F}, P)$ with natural filtration
$(\mathcal{F}_t^L)$.  The L\'evy measure is
$\nu(dz) = c_\gamma|z|^{-1-\gamma}\,dz$, the Poisson random
measure is $N(dt,dz)$, and $\widetilde{N}(dt,dz) =
N(dt,dz) - \nu(dz)\,dt$ is the compensated measure.

Fix $p \in (1,\gamma)$ with conjugate $q = p/(p-1) > 2$.

\begin{definition}[One-parameter objects]\label{def:one-param}
The \emph{one-parameter energy space}, \emph{divergence},
and \emph{derivative} are:
\begin{align}
  H_L &:= L^p([0,T]\times\Omega,\,\mathcal{P},\,
  dt\otimes dP), \label{eq:HL}\\
  \delta_L &: H_L \to L^p(\Omega), \qquad
  \delta_L(u) := \int_0^T u_t\,dL_t, \label{eq:deltaL}\\
  D_L &:= \delta_L^* : L^q(\Omega) \to H_L^*, \qquad
  \langle D_L F, u \rangle = \E[F\cdot\delta_L(u)].
  \label{eq:DL}
\end{align}
The BDG equivalence
$c_p\|u\|_{H_L} \leq \|\delta_L(u)\|_{L^p}
\leq C_p\|u\|_{H_L}$ holds for constants
$c_p, C_p > 0$.
\end{definition}

The factorization~\eqref{eq:factorization} was established
in~\cite{Fontes2026d} on the domain
$\{F \in L^p \cap L^q : F - \E[F] \in
\overline{\im(\delta_L)}^{L^p}\}$.
The key point is that $\overline{\im(\delta_L)}^{L^p}$
is a \emph{proper subspace} of $L^p_0(\Omega) :=
\{F \in L^p : \E[F] = 0\}$ (see
Section~\ref{sec:obstruction}).  The two-parameter extension
enlarges the divergence to achieve surjectivity.


%======================================================================
\section{The Two-Parameter Energy Space}\label{sec:two-param}
%======================================================================

\subsection{Definition}

\begin{definition}[Two-parameter objects]\label{def:two-param}
The \emph{two-parameter energy space} is
\begin{equation}\label{eq:tHL}
  \widetilde{H}_L :=
  L^p\bigl([0,T]\times\R\times\Omega,\;
  \mathcal{P}\otimes\mathcal{B}(\R),\;
  dt\otimes\nu(dz)\otimes dP\bigr),
\end{equation}
equipped with the norm
\begin{equation}\label{eq:tHL-norm}
  \|h\|_{\widetilde{H}_L}
  := \left(\E\!\int_0^T\!\!\int_{\R}
  |h(t,z)|^p\,\nu(dz)\,dt\right)^{1/p}.
\end{equation}
The \emph{two-parameter divergence} is
\begin{equation}\label{eq:tdelta}
  \widetilde{\delta}_L : \widetilde{H}_L \to L^p(\Omega),
  \qquad
  \widetilde{\delta}_L(h) :=
  \int_0^T\!\!\int_{\R} h(t,z)\,\widetilde{N}(dt,dz).
\end{equation}
The \emph{two-parameter derivative} is the adjoint
\begin{equation}\label{eq:tDL}
  \widetilde{D}_L := \widetilde{\delta}_L^* :
  L^q(\Omega) \to \widetilde{H}_L^*,
  \qquad
  \langle \widetilde{D}_L F, h \rangle
  = \E\bigl[F \cdot \widetilde{\delta}_L(h)\bigr]
\end{equation}
for $F \in L^q(\Omega)$,
$h \in \widetilde{H}_L$.
\end{definition}

\begin{remark}[The measure $dt\otimes\nu\otimes dP$]
\label{rem:measure}
The product measure $dt\otimes\nu(dz)\otimes dP$ is
$\sigma$-finite but not finite (since $\nu(\R) = \infty$
for $\gamma < 2$).  The space $\widetilde{H}_L$ is a
well-defined separable Banach space.  The predictability
requirement $\mathcal{P}\otimes\mathcal{B}(\R)$ ensures
that integrands $h(t,z,\omega)$ are predictable in $t$
(adapted to $\mathcal{F}_{t-}^L$) and measurable in $z$.
\end{remark}


\subsection{The linearization map and the $\nu$-singularity obstruction}

The natural candidate for relating the one-parameter and
two-parameter frameworks is the linearization map.  We show
that it does \emph{not} define an embedding.

\begin{definition}[Linearization map]\label{def:linearization}
For $u \in H_L$, define the \emph{linearization}
$\Lambda u : [0,T]\times\R\times\Omega \to \R$ by
\begin{equation}\label{eq:Lambda}
  (\Lambda u)(t,z) := u(t) \cdot z.
\end{equation}
\end{definition}

\begin{proposition}[Properties of $\Lambda$]
\label{prop:Lambda}
\begin{enumerate}[label=(\alph*)]
\item $\Lambda$ does \emph{not} map $H_L$ into
  $\widetilde{H}_L$: the norm
  $\|\Lambda u\|_{\widetilde{H}_L}^p
  = \|u\|_{H_L}^p \cdot \int_{\R}|z|^p\,\nu(dz) = \infty$,
  since $\int|z|^p\,\nu(dz) = \infty$ for $p < \gamma$
  (the integral $\int_0^1 z^{p-1-\gamma}\,dz$ diverges
  when $p < \gamma$).
  The integrand $u(t)\cdot z$ is too singular near $z = 0$
  for the $L^p(\nu)$ norm.
\item Nevertheless, $\Lambda$ intertwines the divergences
  \emph{at the level of stochastic integrals}:
  $\int_0^T\!\int_\R u(t)z\,\widetilde{N}(dt,dz) = \delta_L(u)
  \in L^p(\Omega)$ for $u \in H_L$, even though
  $u(t)z \notin \widetilde{H}_L$.
  The integral converges by cancellation (compensation),
  not by integrability of the integrand.
\end{enumerate}
\end{proposition}

\begin{proof}
(a) $\int_\R|z|^p\nu(dz) = 2c_\gamma\int_0^\infty
z^{p-1-\gamma}\,dz$.  Near $z = 0$: the exponent
$p - 1 - \gamma < -1$ (since $p < \gamma$ and $\gamma < 2$
gives $p - 1 - \gamma < \gamma - 1 - \gamma = -1$), so the
integral diverges.

(b) The identity $\delta_L(u) = \int\!\int u(t)z\,
\widetilde{N}(dt,dz)$ holds because $L_t = \int_0^t\!\int_\R
z\,\widetilde{N}(ds,dz)$ (valid for symmetric stable with
$\gamma > 1$, since $\int_{|z|\leq 1}z\,\nu(dz) = 0$
by symmetry).  The one-parameter BDG inequality
$\|\delta_L(u)\|_{L^p} \leq C_p\|u\|_{H_L}$ (established
in~\cite{Fontes2026d}) shows the integral converges in
$L^p(\Omega)$, even though $u(t)z$ is not in
$L^p(\nu\otimes dt\otimes dP)$.
\end{proof}

\begin{remark}[Two regimes of Poisson integrals]
\label{rem:two-regimes}
The stochastic integrals against $\widetilde{N}(dt,dz)$ fall
into two distinct regimes:
\begin{enumerate}[label=(\roman*)]
\item \textbf{$\nu$-regular integrands:}
  $h \in \widetilde{H}_L = L^p(\nu\otimes dt\otimes dP)$.
  The integral converges because the integrand is
  $p$-th-power integrable against $\nu$.  The defect integrand
  $g(t,z) = \phi(L_{t-}+z) - \phi(L_{t-}) - \phi'(L_{t-})z
  = O(z^2)$ belongs to this regime, since
  $\int_{|z|\leq 1}|z|^{2p}\nu(dz) < \infty$ for
  $p > \gamma/2$ (which holds because $p > 1 > \gamma/2$).
\item \textbf{$\nu$-singular integrands:}
  $h(t,z) = u(t)\cdot z$ with $u \in H_L$.  The integrand
  is $O(z)$ near $z = 0$, too singular for $L^p(\nu)$
  (since $\int|z|^p\nu(dz) = \infty$).  The integral
  converges by cancellation in the compensation
  $\widetilde{N} = N - \nu\,dt$.
\end{enumerate}
The one-parameter $\delta_L$-integrals are $\nu$-singular.
The defect is $\nu$-regular.  A full Clark--Ocone
representation must accommodate both regimes.
\end{remark}


%======================================================================
\section{Two-Parameter BDG Equivalence}\label{sec:BDG}
%======================================================================

The norm equivalence for compensated Poisson integrals
holds for $\nu$-regular integrands.

\begin{theorem}[Two-parameter BDG equivalence]
\label{thm:two-param-BDG}
For $p \in (1,\gamma)$ with $\gamma \in (1,2)$, there exist
constants $\tilde{c}_p, \tilde{C}_p > 0$ (depending on $p$,
$\gamma$, and $T$) such that for all
$h \in \widetilde{H}_L$,
\begin{equation}\label{eq:two-param-BDG}
  \tilde{c}_p\,\|h\|_{\widetilde{H}_L}
  \leq \|\widetilde{\delta}_L(h)\|_{L^p(\Omega)}
  \leq \tilde{C}_p\,\|h\|_{\widetilde{H}_L}.
\end{equation}
\end{theorem}

\begin{proof}[Proof sketch]
The upper bound follows from the Burkholder--Rosenthal
inequality for compensated Poisson integrals
(see~\cite{Dirksen2015}, Theorem~1.1;
or~\cite{MarinelliRockner2014}, Theorem~3.1).
For $p \in (1,2)$ and a predictable integrand
$h \in \widetilde{H}_L$:
\begin{equation}\label{eq:upper-BDG}
  \E\!\left[\left|\int_0^T\!\!\int_{\R}
  h\,\widetilde{N}\right|^p\right]
  \leq C_p\,\E\!\left[\int_0^T\!\!\int_{\R}
  |h(t,z)|^p\,\nu(dz)\,dt\right]
  = C_p\,\|h\|_{\widetilde{H}_L}^p.
\end{equation}
For $p \in (1,2)$, the general Burkholder--Rosenthal
inequality involves both a ``big jump'' term
$\E[\int\!\int|h|^p\,\nu\,dt]$ and a ``small jump'' term
$\E[(\int\!\int|h|^2\,\nu\,dt)^{p/2}]$.
For $p < 2$, the $L^p(\nu)$ term dominates: the individual
large jumps control the $L^p$-behavior of the integral,
and the upper bound simplifies to the
$\widetilde{H}_L$ norm alone.

The lower bound follows from decoupling inequalities
(\cite{delaPenaGine1999}, Chapter~6).  The key step is:
by a tangent sequence argument, the $L^p$ norm of the
compensated Poisson integral is bounded below by
(a constant times) the $L^p$ norm of the corresponding
decoupled sum, which in turn is bounded below by the
$L^p(\nu\otimes dt\otimes dP)$ norm of the integrand.
\end{proof}

\begin{remark}[Scope of the two-parameter BDG]
\label{rem:BDG-scope}
The equivalence~\eqref{eq:two-param-BDG} applies only
to $\nu$-regular integrands ($h \in \widetilde{H}_L$).
It does \emph{not} apply to $\nu$-singular integrands
such as $h(t,z) = u(t)z$ (which define $\delta_L$-integrals);
for these, the one-parameter BDG
$\|\delta_L(u)\|_{L^p} \asymp \|u\|_{H_L}$
is the correct estimate (see~\cite{Fontes2026d}).
The two BDG inequalities govern different classes of
integrands and are \emph{not} special cases of each other.
\end{remark}

\begin{corollary}[Recovery map for $\nu$-regular integrands]
\label{cor:recovery}
The two-parameter divergence
$\widetilde{\delta}_L : \widetilde{H}_L \to
\overline{\im(\widetilde{\delta}_L)}^{L^p}$
is a Banach space isomorphism onto its range.  The
\emph{two-parameter recovery map}
$\widetilde{\iota}_L^{-1} :
\im(\widetilde{\delta}_L)
\to \widetilde{H}_L$ is the bounded inverse
constructed from the BDG equivalence, with operator norm
$\|\widetilde{\iota}_L^{-1}\| \leq 1/\tilde{c}_p$.
\end{corollary}


%======================================================================
\section{The Two-Parameter Factorization and Clark--Ocone}
\label{sec:factorization}
%======================================================================

\subsection{Surjectivity}

The critical difference from the one-parameter case is that the
two-parameter divergence is surjective onto $L^p_0(\Omega)$.

\begin{theorem}[Surjectivity of $\widetilde{\delta}_L$]
\label{thm:surjectivity}
The closure $\overline{\im(\widetilde{\delta}_L)}^{L^p}
= L^p_0(\Omega)$.  That is, every centered
$F \in L^p(\Omega)$ lies in the $L^p$-closure of the range
of $\widetilde{\delta}_L$.
\end{theorem}

\begin{proof}
We show that the annihilator
$\im(\widetilde{\delta}_L)^{\perp} = \{0\}$ in
$L^q(\Omega)$ (here $\perp$ denotes the annihilator
under the $L^p$--$L^q$ duality).

Suppose $G \in L^q(\Omega)$ satisfies
$\E[G \cdot \widetilde{\delta}_L(h)] = 0$ for all
$h \in \widetilde{H}_L$.  By definition of
$\widetilde{D}_L$, this means
$\langle \widetilde{D}_L G, h \rangle = 0$ for all $h$,
hence $\widetilde{D}_L G = 0$ in $\widetilde{H}_L^*$.

We claim that $\widetilde{D}_L G = 0$ implies
$G = \E[G]$.  To see this, note that
$\langle \widetilde{D}_L G, h \rangle
= \E[G \cdot \widetilde{\delta}_L(h)] = 0$ for all
predictable $h \in \widetilde{H}_L$.
In particular, for $h = \mathbf{1}_{(a,b]}\mathbf{1}_A
\cdot g(z)$ with $A \in \mathcal{F}_a^L$ and
$g \in L^p(\R, \nu)$:
\[
  \E\!\left[\mathbf{1}_A\,G \cdot \int_a^b\!\!\int_\R
  g(z)\,\widetilde{N}(ds,dz)\right] = 0.
\]
Since this holds for all $A \in \mathcal{F}_a^L$,
it follows that
$\E[G\cdot\int_a^b\!\int g(z)\widetilde{N}(ds,dz)
\mid\mathcal{F}_a^L] = 0$ a.s.
The compensated Poisson integrals
$\int_a^b\!\int g(z)\widetilde{N}(ds,dz)$ for varying
$a < b$ and $g \in L^p(\R,\nu)$ detect both the times
and sizes of all jumps of $L$, generating
$\mathcal{F}_T^L$.  The conditional orthogonality
for all $A \in \mathcal{F}_a^L$ then forces
$G = \E[G]$ a.s.

Therefore $\im(\widetilde{\delta}_L)^{\perp}
\cap L^q_0(\Omega) = \{0\}$, and the Hahn--Banach
theorem gives $\overline{\im(\widetilde{\delta}_L)}^{L^p}
\supset L^p_0(\Omega)$.  The reverse inclusion is
immediate from $\E[\widetilde{\delta}_L(h)] = 0$.
\end{proof}

\begin{remark}[Why the one-parameter divergence fails]
\label{rem:one-param-fails}
The analogous argument for $\delta_L$ would require
that $\E[G \cdot \delta_L(u)] = 0$ for all
$u \in H_L$ implies $G = \E[G]$.  This fails.
Let $G_0 = \sum_{s\leq T}(\Delta L_s)^2
\mathbf{1}_{|\Delta L_s|>1}$
(a compound Poisson sum, hence $G_0 \in L^r$ for all $r$)
and $G = G_0 - \E[G_0]$.  Then $G$ is centered and nonzero.
For any $u \in H_L$, $\delta_L(u) = \int u_t\,dL_t$ is
\emph{odd} in the jump sizes (linear in $\Delta L_s$),
while $G$ is \emph{even} (depends on $(\Delta L_s)^2$).
By symmetry of the L\'evy measure $\nu$,
$\E[G\cdot\delta_L(u)] = 0$, so $D_L G = 0$ but
$G \neq 0$.
\end{remark}


\subsection{The factorization theorem}

\begin{theorem}[Two-parameter factorization]
\label{thm:two-param-factorization}
For $p \in (1,\gamma)$ with $\gamma \in (1,2)$, the identity
\begin{equation}\label{eq:two-param-fact}
  F - \E[F]
  = \widetilde{\delta}_L\bigl(
  \widetilde{\iota}_L^{-1}(\widetilde{D}_L F)\bigr)
\end{equation}
holds for all $F \in L^p(\Omega) \cap L^q(\Omega)$.
\end{theorem}

\begin{proof}
By Theorem~\ref{thm:surjectivity}, $F - \E[F] \in L^p_0
= \overline{\im(\widetilde{\delta}_L)}^{L^p}$.  By the
BDG equivalence (Theorem~\ref{thm:two-param-BDG}), there
exists a unique $h_F \in \widetilde{H}_L$ such that
$\widetilde{\delta}_L(h_F) = F - \E[F]$.  By the BDG lower
bound, $h_F = \widetilde{\iota}_L^{-1}(F - \E[F])$.

It remains to identify
$h_F = \widetilde{\iota}_L^{-1}(\widetilde{D}_L F)$.
For any $g \in \widetilde{H}_L$:
\begin{align*}
  \langle \widetilde{D}_L F, g \rangle
  &= \E[F \cdot \widetilde{\delta}_L(g)]
  = \E[(F - \E[F]) \cdot \widetilde{\delta}_L(g)]\\
  &= \E[\widetilde{\delta}_L(h_F) \cdot
  \widetilde{\delta}_L(g)].
\end{align*}
This identifies $\widetilde{D}_L F$ as the functional
$g \mapsto \E[\widetilde{\delta}_L(h_F)\cdot
\widetilde{\delta}_L(g)]$, which is the image of $h_F$
under the map $\widetilde{H}_L \to \widetilde{H}_L^*$
induced by the bilinear form
$(h, g) \mapsto \E[\widetilde{\delta}_L(h)\cdot
\widetilde{\delta}_L(g)]$.  The recovery map
$\widetilde{\iota}_L^{-1}$ inverts this identification.
\end{proof}


\subsection{The Clark--Ocone formula}

\begin{theorem}[$L^p$ Clark--Ocone for stable processes]
\label{thm:clark-ocone}
For $p \in (1,\gamma)$ and $F \in L^p(\Omega) \cap L^q(\Omega)$,
there exists a unique $h_F \in \widetilde{H}_L$ such that
\begin{equation}\label{eq:clark-ocone}
  F - \E[F]
  = \int_0^T\!\!\int_{\R} h_F(t,z)\,\widetilde{N}(dt,dz),
\end{equation}
with the estimate
\begin{equation}\label{eq:CO-estimate}
  \tilde{c}_p\,\|h_F\|_{\widetilde{H}_L}
  \leq \|F - \E[F]\|_{L^p(\Omega)}
  \leq \tilde{C}_p\,\|h_F\|_{\widetilde{H}_L}.
\end{equation}
The integrand $h_F$ is predictable (in $t$) and
$\nu$-measurable (in $z$).
\end{theorem}

\begin{proof}
This is immediate from the surjectivity
(Theorem~\ref{thm:surjectivity}), the BDG equivalence
(Theorem~\ref{thm:two-param-BDG}), and the factorization
(Theorem~\ref{thm:two-param-factorization}).
\end{proof}

\begin{remark}[Comparison with existing results]
\label{rem:comparison}
The $L^2$ Clark--Ocone formula for L\'evy processes was
established by L{\o}kka~\cite{Lokka2004} and
Di~Nunno--{\O}ksendal--Proske~\cite{DiNunnoOksendalProske2009}
using the Picard difference operator $D_{s,z}F =
F(\omega + z\epsilon_s) - F(\omega)$ and chaos expansions.
Their formula reads
$F - \E[F] = \int\!\int \E[D_{s,z}F|\mathcal{F}_{s-}]
\,\widetilde{N}(ds,dz)$.

Theorem~\ref{thm:clark-ocone} extends this to $L^p$ with
$p < 2$ via operator methods, without chaos expansions or
difference operators.  The integrand
$h_F = \widetilde{\iota}_L^{-1}(\widetilde{D}_L F)$ is
defined purely through the adjoint structure.  The
connection to the Picard operator is discussed in
Section~\ref{sec:discussion}.
\end{remark}


%======================================================================
\section{Two-Regime Decomposition}\label{sec:decomposition}
%======================================================================

We now analyze the structure of the two-parameter integrand
$h_F$ for functionals of the form $F = \phi(L_T)$, revealing
the relationship between the one-parameter It\^o formula and
the two-parameter Clark--Ocone representation.

\subsection{The two-regime decomposition}

The Banach It\^o formula~\eqref{eq:banach-ito} writes
$F - \E[F] = \delta_L(\phi'(L_{\cdot-}))
+ [\Gamma_L^{\phi}(T) - \E[\Gamma_L^{\phi}(T)]]$.
Each summand is a centered $L^p$ random variable,
hence (by Theorem~\ref{thm:clark-ocone}) has a unique
$\widetilde{H}_L$-representative.  However, these
representatives have fundamentally different character.

\begin{theorem}[Two-regime decomposition]\label{thm:defect-decomposition}
Let $\phi \in C_b^3(\R)$ and $F = \phi(L_T)$.  Then:
\begin{enumerate}[label=(\alph*)]
\item The defect $\Gamma_L^\phi(T)$ decomposes as
  $\Gamma_L^\phi(T) = M_T^{(1)} + A_T^{(1)}$, where
  the martingale part is
  \begin{equation}\label{eq:defect-martingale}
    M_T^{(1)} := \int_0^T\!\!\int_{\R}
    g_\phi(s,z)\,\widetilde{N}(ds,dz)
    = \widetilde{\delta}_L(g_\phi),
  \end{equation}
  with $g_\phi(s,z) :=
  \phi(L_{s-}+z) - \phi(L_{s-}) - \phi'(L_{s-})z$,
  and the compensator is
  $A_T^{(1)} = \int_0^T\mathcal{L}\phi(L_s)\,ds$.
  The integrand $g_\phi$ lies in $\widetilde{H}_L$:
  it satisfies $|g_\phi(s,z)| \leq
  \tfrac{1}{2}\|\phi''\|_\infty z^2$ for $|z|\leq 1$
  and $\int_{|z|\leq 1}|z|^{2p}\nu(dz) < \infty$
  (since $p > 1 > \gamma/2$).
\item The stochastic integral
  $\delta_L(\phi'(L_{\cdot-}))$ also has a unique
  $\widetilde{H}_L$-representative $h^{(0)} \in
  \widetilde{H}_L$ with $\widetilde{\delta}_L(h^{(0)})
  = \delta_L(\phi'(L_{\cdot-}))$, but this representative
  is \emph{not} the function $\phi'(L_{t-})\cdot z$
  (which lies outside $\widetilde{H}_L$).
  Rather, $h^{(0)}$ is an abstract element whose
  $\widetilde{N}$-integral happens to equal
  the $dL$-integral of $\phi'(L_{\cdot-})$.
\item The full Clark--Ocone integrand satisfies
  $h_F = h^{(0)} + h^{(1)}$ where $h^{(1)}$ is the
  $\widetilde{H}_L$-representative of the centered defect.
\end{enumerate}
\end{theorem}

\begin{proof}
(a) By Taylor's theorem, $|g_\phi(s,z)| \leq
\tfrac{1}{2}\|\phi''\|_\infty z^2$ for $|z| \leq 1$
and $|g_\phi(s,z)| \leq 2\|\phi\|_\infty +
\|\phi'\|_\infty|z|$ for $|z| > 1$.  Therefore
\[
  \|g_\phi\|_{\widetilde{H}_L}^p
  \leq C\,\E\!\int_0^T\!\left[
  \int_{|z|\leq 1}|z|^{2p}\nu(dz)
  + \int_{|z|>1}(1+|z|)^p\nu(dz)\right]dt.
\]
The first integral:
$\int_{|z|\leq 1}|z|^{2p}\nu(dz)
= 2c_\gamma\int_0^1 z^{2p-1-\gamma}\,dz < \infty$
iff $2p > \gamma$, i.e., $p > \gamma/2$.
Since $p > 1 > \gamma/2$ (for $\gamma < 2$), this holds.
The second integral converges since $p < \gamma$.
Thus $g_\phi \in \widetilde{H}_L$ and
$M_T^{(1)} = \widetilde{\delta}_L(g_\phi)$.

The Doob decomposition gives
$\Gamma_L^\phi(T) = M_T^{(1)} + A_T^{(1)}$ where
$A_T^{(1)} = \int_0^T\!\int_\R g_\phi(s,z)\nu(dz)\,ds
= \int_0^T\mathcal{L}\phi(L_s)\,ds$
(with $\mathcal{L} = -(-\Delta)^{\gamma/2}$, using
symmetry of $\nu$ to identify the truncated and
untruncated generators).
The centered defect is therefore
$\Gamma_L^\phi(T) - \E[\Gamma_L^\phi(T)]
= M_T^{(1)} + (A_T^{(1)} - \E[A_T^{(1)}])$.
By Theorem~\ref{thm:clark-ocone}, each centered summand
has a unique $\widetilde{H}_L$-representative:
$g_\phi$ for $M_T^{(1)}$, and an abstract element
for the compensator correction $A_T^{(1)} - \E[A_T^{(1)}]$.

(b) $\delta_L(\phi'(L_{\cdot-})) \in L^p_0(\Omega)$
(centered by the BDG construction).  By
Theorem~\ref{thm:clark-ocone}, it has a unique
$\widetilde{H}_L$-representative $h^{(0)}$.
The natural integrand $\phi'(L_{t-})z$ satisfies
$\int|z|^p\nu(dz) = \infty$ (Proposition~\ref{prop:Lambda}(a)),
so $\phi'(L_{t-})z \notin \widetilde{H}_L$.
The representative $h^{(0)}$ exists by the abstract
surjectivity theorem but does not admit a simple formula
in terms of $\phi'$ and $z$.

(c) By uniqueness (Theorem~\ref{thm:clark-ocone})
and linearity of $\widetilde{\delta}_L$.
\end{proof}

\begin{remark}[The two regimes, revisited]
\label{rem:regimes-revisited}
The decomposition reveals a fundamental asymmetry:
\begin{enumerate}[label=(\roman*)]
\item The \emph{defect part} has a natural, explicit
  $\widetilde{H}_L$-integrand $g_\phi(s,z) = O(z^2)$.
  The quadratic vanishing at $z = 0$ puts it squarely
  in $\widetilde{H}_L$.
\item The \emph{stochastic integral part} has a
  natural integrand $\phi'(L_{t-})z = O(z)$ that is
  too singular for $\widetilde{H}_L$.  Its
  $\widetilde{H}_L$-representative exists abstractly
  but is non-explicit.
\end{enumerate}
This is the precise sense in which the Leibniz defect
is ``better behaved'' than the stochastic integral in
the two-parameter framework: the defect is $\nu$-regular,
while the stochastic integral is $\nu$-singular.

In the Hilbert case (Brownian motion), the L\'evy measure
$\nu = 0$, there are no jumps, and the two-parameter
space is trivial.  The passage from stable ($\gamma < 2$)
to Brownian ($\gamma = 2$) is not a continuous limit
in the L\'evy measure topology: it is a singular transition
in which the jump structure vanishes entirely and the
defect $\Gamma_L^\phi \to 0$.
\end{remark}

\begin{remark}[Defect norm and regularity]
\label{rem:distance}
Although $\Lambda(H_L) \not\subset \widetilde{H}_L$
(the linear-in-$z$ functions are too singular), the
\emph{defect norm}
$\|g_\phi\|_{\widetilde{H}_L}$ measures
the size of the chain-rule correction.
Since the leading behavior $g_\phi(s,z) \approx
\tfrac{1}{2}\phi''(L_{s-})z^2$ is determined by
$\phi''$, the defect norm is controlled by
$\|\phi''\|_\infty$ (Proposition~\ref{prop:regularity}):
larger curvature of $\phi$ produces a larger defect.
Higher-order smoothness ($\phi \in C_b^k$ for $k > 3$)
does not improve the norm, as the leading $z^2$ term
dominates near $z = 0$.
\end{remark}


%======================================================================
\section{Regularity Implications}\label{sec:regularity}
%======================================================================

The defect integrand $g_\phi$ encodes the regularity of
$\phi$ through its $z$-decay near $z = 0$.  Since $g_\phi$
lies in $\widetilde{H}_L$, its norm provides a quantitative
measure of the chain-rule correction.

\begin{proposition}[Defect norm controls regularity]
\label{prop:regularity}
Let $\phi \in C_b^3(\R)$.  The defect integrand
$g_\phi(s,z) := \phi(L_{s-}+z) - \phi(L_{s-})
- \phi'(L_{s-})z$ satisfies
\begin{equation}\label{eq:defect-norm-bound}
  \|g_\phi\|_{\widetilde{H}_L}^p
  \leq C\!\left(\|\phi''\|_\infty^p
  \int_{|z|\leq 1}|z|^{2p}\nu(dz)
  + (2\|\phi\|_\infty + \|\phi'\|_\infty)^p
  \int_{|z|>1}(1+|z|)^p\nu(dz)\right)\!T,
\end{equation}
where both integrals are finite:
$\int_{|z|\leq 1}|z|^{2p}\nu(dz) < \infty$ because
$2p > \gamma$ (since $p > 1 > \gamma/2$),
and $\int_{|z|>1}(1+|z|)^p\nu(dz) < \infty$ because
$p < \gamma$.

Consequently, the defect norm is controlled by the
$C_b^2$ norm of $\phi$:
$\|g_\phi\|_{\widetilde{H}_L}
\leq C(\gamma,p,T)\,
(\|\phi''\|_\infty + \|\phi\|_\infty + \|\phi'\|_\infty)$.
\end{proposition}

\begin{proof}
By Taylor's theorem with remainder:
$|g_\phi(s,z)| \leq \tfrac{1}{2}\|\phi''\|_\infty\,z^2$
for $|z| \leq 1$ (using $\phi \in C_b^3$).
For $|z| > 1$:
$|g_\phi(s,z)| \leq |\phi(L_{s-}+z)| + |\phi(L_{s-})|
+ |\phi'(L_{s-})||z|
\leq 2\|\phi\|_\infty + \|\phi'\|_\infty|z|$.
Taking $p$-th powers and integrating against
$\nu(dz)\,ds\,dP$ gives~\eqref{eq:defect-norm-bound}.
\end{proof}

\begin{remark}[The integrability threshold]
\label{rem:threshold}
The condition $2p > \gamma$ (equivalently $p > \gamma/2$)
is the threshold for $g_\phi \in \widetilde{H}_L$.
Since $\gamma \in (1,2)$ and $p > 1$, this is always
satisfied: $p > 1 > \gamma/2$.
\end{remark}

\begin{remark}[Connection to the fractional Laplacian]
\label{rem:besov}
The compensator $A_T^{(1)} = \int_0^T\!\int_\R
g_\phi(s,z)\,\nu(dz)\,ds
= \int_0^T\mathcal{L}\phi(L_s)\,ds$, where
$\mathcal{L} = -(-\Delta)^{\gamma/2}$, connects the
defect norm to the regularity of $\phi$ under the
fractional Laplacian.  Specifically,
$\|\mathcal{L}\phi\|_\infty < \infty$ iff $\phi$ is in
the domain of $(-\Delta)^{\gamma/2}$ as a bounded
operator, which for $\phi \in C_b^2$ holds when
$\gamma < 2$.  The defect norm
$\|g_\phi\|_{\widetilde{H}_L}$ refines this to an
$L^p$-stochastic setting: it measures not just the
pointwise action of $\mathcal{L}$ but the $p$-th
moment of the jump-by-jump correction along the
path of $L$.
\end{remark}


%======================================================================
\section{Discussion}\label{sec:discussion}
%======================================================================

\subsection{The complete picture}

The operator factorization program now spans three levels:

\medskip\noindent\textbf{Level 1: Hilbert.}
Energy space $H = L^2(dt\otimes dP)$.  The factorization
$(\Id - \E) = \delta \circ D$ (with Riesz identification)
gives Clark--Ocone.  The It\^o formula has correction
$\frac{1}{2}\|\Pi DF\|_H^2$.  No defect ($\Gamma = 0$).

\medskip\noindent\textbf{Level 2: Banach, one-parameter.}
Energy space $H_L = L^p(dt\otimes dP)$.  The factorization
holds on $\overline{\im(\delta_L)}^{L^p} \subsetneq L^p_0$.
The It\^o formula has the Leibniz defect $\Gamma_L^{\phi}$
as correction.  No Clark--Ocone for general functionals.

\medskip\noindent\textbf{Level 3: Banach, two-parameter.}
Energy space $\widetilde{H}_L = L^p(dt\otimes\nu\otimes dP)$.
The factorization holds on all of $L^p_0$.  Clark--Ocone
holds for all centered $F \in L^p \cap L^q$.  The defect
integrand $g_\phi = O(z^2)$ is $\nu$-regular and has an
explicit $\widetilde{H}_L$-representative, while the
$dL$-integral $\delta_L(\phi'(L_{\cdot-}))$ is $\nu$-singular
and has only an abstract representative.

\medskip
The passage from Level~1 to Level~2 introduced the defect
as a correction term.  The passage from Level~2 to Level~3
reinterprets the defect as an obstruction and resolves it
by enlarging the energy space.

\subsection{Relationship to the Picard operator}

In $L^2$, the Clark--Ocone formula of
L{\o}kka~\cite{Lokka2004} uses the Picard difference operator
$D_{s,z}F = F(\omega + z\epsilon_s) - F(\omega)$, giving
\[
  F - \E[F] = \int_0^T\!\!\int_\R
  \E[D_{s,z}F \mid \mathcal{F}_{s-}]\,
  \widetilde{N}(ds,dz).
\]
Our Theorem~\ref{thm:clark-ocone} gives the same structural
formula with $h_F$ replacing $\E[D_{s,z}F|\mathcal{F}_{s-}]$.
In $L^2$, the two must agree (by uniqueness of the representation);
thus the operator-theoretic integrand
$\widetilde{\iota}_L^{-1}(\widetilde{D}_L F)$ coincides
with $\E[D_{s,z}F|\mathcal{F}_{s-}]$ in $L^2$.

For $p < 2$, the relationship is more subtle.  The Picard
operator is defined pathwise (as a difference), while our
$\widetilde{D}_L$ is defined as an adjoint (through duality).
Whether the two agree in $L^p$ for $p < 2$---and whether the
Picard operator extends to a bounded operator in the
$L^p$ setting---is an open question.

The operator-theoretic approach has the advantage of not
requiring pathwise constructions or chaos expansions: the
integrand exists by abstract Banach space theory (BDG +
surjectivity), regardless of whether it admits a
pointwise formula.


\subsection{Open problems}

\begin{enumerate}[label=(\roman*)]
\item \textbf{Explicit computation of
$\widetilde{\iota}_L^{-1}$.}
The recovery map is constructed abstractly from BDG.
For $L^2$, it reduces to the conditional expectation
$h_F(t,z) = \E[D_{t,z}F|\mathcal{F}_{t-}]$.  For
$L^p$ with $p < 2$: is there a conditional-expectation-based
formula, and if so, what replaces the Picard operator $D_{t,z}$
inside the conditional expectation?

\item \textbf{The duality map and nonlinear projections.}
The $L^p$--$L^q$ duality map
$J_q(g) = |g|^{q-2}g/\|g\|^{q-2}$ is the natural
candidate to replace Riesz identification.  Does $J_q$ interact
with the predictable projection and the jump structure in a way
that produces a computable integrand?

\item \textbf{Refined regularity scales.}
The defect norm $\|g_\phi\|_{\widetilde{H}_L}$ is a
coarse measure of regularity.  Can one define finer
regularity scales based on the $z$-decay of $g_F$
(analogous to Besov spaces $B^\alpha_{p,q}$) that
characterize the domain of the fractional Laplacian
in probabilistic terms?

\item \textbf{Extension to $\gamma \in (0,1]$.}
For $\gamma \leq 1$, the range $p \in (1,\gamma)$ is empty,
so the entire framework requires reformulation---either
working with $p \in (0,1)$ (quasi-Banach spaces) or with
$p > \gamma$ (which changes the BDG landscape, since the
process has finite variation and different moment behavior).
The two-parameter space
$\widetilde{H}_L = L^p(dt\otimes\nu\otimes dP)$ should
still be natural, but the surjectivity argument and
BDG equivalence need new ingredients.

\item \textbf{Non-stable jump processes.}
For general L\'evy processes with L\'evy measure $\nu$
(not necessarily stable), the two-parameter framework
should extend with $\widetilde{H}_L =
L^p(dt\otimes\nu\otimes dP)$.  The BDG equivalence
would depend on the tail behavior of $\nu$.

\item \textbf{Applications to hedging.}
In financial models driven by stable processes, the
one-parameter integrand $\phi'(L_{t-})$ is the delta hedge,
and the defect is the hedging error from jumps.  The
two-parameter integrand $h_F(t,z)$ provides the
\emph{jump-size-dependent hedge}: for each potential jump
of size $z$, the required position adjustment.  Making
this precise in the context of incomplete markets is an
open problem.
\end{enumerate}


%======================================================================
\section*{Acknowledgments}
%======================================================================
AI tools were used as interactive assistants for
mathematical exploration and drafting.  The author is
responsible for all content.


%======================================================================
\begin{thebibliography}{10}

\bibitem{Applebaum2009}
D.~Applebaum.
\newblock \emph{L\'evy Processes and Stochastic Calculus}.
\newblock Cambridge University Press, 2nd edition, 2009.

\bibitem{delaPenaGine1999}
V.~de~la~Pe\~na and E.~Gin\'e.
\newblock \emph{Decoupling: From Dependence to Independence}.
\newblock Springer, 1999.

\bibitem{DiNunnoOksendalProske2009}
G.~Di~Nunno, B.~{\O}ksendal, and F.~Proske.
\newblock \emph{Malliavin Calculus for L\'evy Processes with
Applications to Finance}.
\newblock Springer, 2009.

\bibitem{Dirksen2015}
S.~Dirksen.
\newblock Itô isomorphisms for $L^p$-valued Poisson stochastic
integrals.
\newblock \emph{Ann.\ Probab.}, 42(6):2595--2643, 2014.

\bibitem{Fontes2026a}
R.~Fontes.
\newblock Operator factorization and the Clark--Ocone formula
in stochastic calculus.
\newblock arXiv preprint, 2026.

\bibitem{Fontes2026b}
R.~Fontes.
\newblock Extension to rough fractional Brownian motion.
\newblock arXiv preprint, 2026.

\bibitem{Fontes2026c}
R.~Fontes.
\newblock An It\^o formula via product rules in Hilbert energy
spaces.
\newblock arXiv preprint, 2026.

\bibitem{Fontes2026d}
R.~Fontes.
\newblock Operator factorization in Banach energy spaces:
Leibniz defects for stable processes.
\newblock arXiv preprint, 2026.

\bibitem{Fontes2026e}
R.~Fontes.
\newblock PDE generators from energy space structure.
\newblock arXiv preprint, 2026.

\bibitem{Fontes2026f}
R.~Fontes.
\newblock An It\^o formula via Leibniz defects in Banach energy
spaces and PDE generators for stable processes.
\newblock arXiv preprint, 2026.

\bibitem{Lokka2004}
A.~L{\o}kka.
\newblock Martingale representation of functionals of L\'evy
processes.
\newblock \emph{Stoch.\ Anal.\ Appl.}, 22(4):867--892, 2004.

\bibitem{MarinelliRockner2014}
C.~Marinelli and M.~R\"ockner.
\newblock On the maximal inequalities of Burkholder, Davis and
Gundy.
\newblock \emph{Expo.\ Math.}, 34(1):1--26, 2016.

\bibitem{Sato1999}
K.-i.~Sato.
\newblock \emph{L\'evy Processes and Infinitely Divisible
Distributions}.
\newblock Cambridge University Press, 1999.

\end{thebibliography}

\end{document}
